\documentclass[12pt]{article}

\usepackage[utf8]{inputenc}
\usepackage{amsmath, amssymb, amsthm}
\usepackage{geometry}
\usepackage{graphicx}
\usepackage{hyperref}
\usepackage{parskip}
\usepackage{booktabs}
\usepackage{float}

\geometry{a4paper, margin=1in}

\title{Warped Hybrid Slice Sampler (WHSS): Gradient-Free MCMC for Highly Constrained, Anisotropic Spaces}
\author{}
\date{}

\begin{document}

\maketitle

\section*{Motivation}
Current state-of-the-art MCMC algorithms, such as the No-U-Turn Sampler (NUTS), rely heavily on the gradient information of the target density to glide through high-dimensional, non-linear, and complex target spaces. However, in many real-world scenarios, the target density itself is strictly non-differentiable or mathematically opaque (black-box). This limits the applicability of Hamiltonian-based algorithms. Furthermore, when these non-differentiable targets are bounded by strict geometric constraints (polytopes) and exhibit extreme anisotropy (skewness), traditional algorithms struggle to maintain efficiency. This presents a significant gap in the literature. Two prominent real-world examples of this paradigm are:
\begin{enumerate}
    \item \textbf{Metabolic Flux Analysis (MFA):} In computational biology, simulating genome-scale metabolic networks requires uniformly sampling the steady-state flux cone. This space is defined by the null space of a stoichiometric matrix ($Sv = 0$) bounded by strict thermodynamic constraints ($lb \le v \le ub$). The resulting polytope is highly skewed and restricted to hundreds of dimensions, making it difficult to navigate.
    \item \textbf{Quantitative Finance (CVaR Allocation):} Institutional portfolio optimization often targets Historical Conditional Value at Risk (CVaR), a strictly non-differentiable risk metric based on empirical tail losses. The portfolio weights are bounded by the standard financial simplex ($\sum_{i=1}^n w_i = 1$) alongside strict regulatory position caps ($w_i \le c$).
\end{enumerate}

\section*{Proposed Algorithm: Warped Hybrid Slice Sampler}
To overcome this gap, we propose the Warped Hybrid Slice Sampler (WHSS). Instead of a standard random walk, our method introduces a novel mechanism to empirically estimate the geometric skew of a constrained space without requiring gradient calculations. The algorithm operates via the following key innovations:

\begin{itemize}
    \item \textbf{vMF Spherical Warm-up:} We initiate exploration using a von Mises-Fisher (vMF) directional distribution. Unlike standard Hit-and-Run which gets trapped in sharp corners, vMF proposes un-skewed directional vectors that bounce across the interior volume of the polytope, identifying a diverse set of "structural anchors."
    \item \textbf{Empirical Space Warping (The Core Novelty):} From these anchors, we construct a robust empirical covariance matrix $V$. We extract its Cholesky decomposition, $L = (V V^T)^{1/2}$. We use $L^{-1}$ as a global transformation matrix to warp the target space. This mathematically eliminates the anisotropic skew, turning long, narrow polytope constraints into well-conditioned, isotropic spheres. 
    \item \textbf{Hybrid Slice Sampling:} In this newly isotropic space, we deploy a multi-chain slice sampler. Because the narrow constraints have been "inflated" by the $L$-matrix, the slice sampler can take large independent steps. This reduces Markov Chain autocorrelation (ESS) without computing a gradient.
\end{itemize}

To clarify our positioning, WHSS does not claim execution speed (wall-clock time) advantages over simpler methods, as gradient-free geometric matrix calculations carry overhead. Instead, we focus on \textit{mathematical efficiency and sample quality}. Therefore, we evaluate our algorithm using \textbf{Algorithmic Efficiency (ESS / 1k NFE)}, which measures: "Per 1,000 algorithmic density/constraint evaluations, how many purely independent, decorrelated samples do we generate?"

\section{Experimental Configurations and Metric Rigor}
To ensure absolute mathematical fairness and transparency, we explicitly define the computational budgets (Number of Function Evaluations, NFE) and burn-in parameters for all experiments. The WHSS algorithm utilizes a one-time \emph{Geometric Warm-Up} to fit its affine warp matrix. These target evaluations are strictly discarded before the MCMC phase begins. Furthermore, depending on the test objective, MCMC chains are either sliced to remove transient drift (for pure steady-state metrics) or evaluated from a cold-start to penalize geometric trapping.

\begin{itemize}
    \item \textbf{Dimensional Scaling:} Total Budget $N=25{,}000$ per algorithm. WHSS Geometric Warm-Up: $1000 \times D$ evaluations. \emph{MCMC Burn-in:} 20\% ($5{,}000$ samples) strictly discarded for all algorithms before computing ESS to eliminate transient drift.
    \item \textbf{Mean Squared Jump Distance (MSJD):} Total Budget $N=15{,}000$ per algorithm. WHSS Geometric Warm-Up: $2{,}500$ evaluations. \emph{MCMC Burn-in:} 0\% (None). The metric relies on the average jump distance between all consecutive points, capturing both transient and stationary jumps symmetrically.
    \item \textbf{Autocorrelation (ACF):} Total Budget $N=25{,}000$ per algorithm. WHSS Geometric Warm-Up: $2{,}500$ evaluations. \emph{MCMC Burn-in:} 0\% (None). This maps exactly how fast an algorithm physically forgets its cold-starting point.
    \item \textbf{Quantitative Finance (CVaR):} Total Budget $N=60{,}000$. WHSS Geometric Warm-Up: $1{,}500$ evaluations. \emph{MCMC Burn-in:} 0\% (None). Tests the full-chain mixing capability starting from the naive equally-weighted portfolio ($w_i = 1/D$).
    \item \textbf{Biological Flux (Core 24D Uniform):} Total Budget $N=500{,}000$. WHSS Geometric Warm-Up: $50{,}000$ evaluations. \emph{MCMC Burn-in:} 50\% ($250{,}000$ samples) strictly discarded for all algorithms before ESS calculation to ensure pure steady-state metrics.
    \item \textbf{Biological Flux (Core 24D Non-Uniform):} Total Budget $N=25{,}000$. WHSS Geometric Warm-Up: $5{,}000$ evaluations. \emph{MCMC Burn-in:} 50\% ($12{,}500$ samples) strictly discarded for all algorithms prior to ESS calculation.
    \item \textbf{Biological Flux (Genome-Scale 542D Uniform):} Total Budget $N=150{,}000$. WHSS Geometric Warm-Up: $15{,}000$ evaluations. \emph{MCMC Burn-in:} 50\% ($75{,}000$ samples) strictly discarded.
\end{itemize}

\section{Mixing Efficiency and Autocorrelation}
\textbf{Motivation:} In highly skewed spaces, isotropic algorithms take small steps, leading to highly correlated samples that fail to map the true distribution. Autocorrelation measures this failure. We included the Hit-and-Run (HRSS) baseline (the legacy standard for bounded spaces) and the Dikin Walk (a mathematically sophisticated barrier-method that uses the log-barrier Hessian to navigate constraints). 

\textbf{Problem Setup:} We defined a $D=30$ dimensional Gaussian target density $f(x) \propto \exp(-\frac{1}{2} x^T \Sigma^{-1} x)$, where the covariance matrix $\Sigma$ is engineered to have a condition number ($\kappa = 1000$). The space is bounded by a strict hypercube constraint $\mathcal{P} = \{x \in \mathbb{R}^{30} \mid -1 \le x_i \le 1\}$. The experiment was run across 10 parallel chains for 25,000 samples per chain.

\textbf{Graph Interpretation:} 
The x-axis represents the Lag ($k$), which is the number of steps between two samples in the Markov Chain. The y-axis represents the Autocorrelation Function (ACF) $\rho_k$. A curve that stays near 1.0 means the algorithm is taking small steps and is mixing poorly. A curve that drops rapidly to 0.0 means the algorithm is successfully taking independent steps across the target space.

\textbf{Observations:}
\begin{itemize}
    \item \textbf{Dikin Walk} struggled to mix. Its autocorrelation remained at exactly $1.000$ out to 250 lags, indicating it was geometrically trapped in a narrow corner.
    \item \textbf{HRSS} exhibited slow mixing due to the anisotropy, retaining an autocorrelation of $0.724$ even at lag 250.
    \item \textbf{WHSS}, utilizing its adaptive $L$-matrix preconditioner, successfully neutralized the geometric skew. Its autocorrelation decreased to $0.086$ at lag 250, proving it successfully decorrelates while baselines struggle.
\end{itemize}

\begin{figure}[H]
    \centering
    \includegraphics[width=0.8\textwidth]{acf_plot.png}
    \caption{Autocorrelation vs Lag in a 30D Skewed Polytope. WHSS successfully decorrelates while Dikin Walk completely freezes.}
\end{figure}

\section{Geometric Mobility (Mean Squared Jump Distance)}
\textbf{Motivation:} Autocorrelation is a statistical metric, but Mean Squared Jump Distance (MSJD) is a pure physical metric. It measures the raw Euclidean distance covered by the algorithm's accepted proposals. In tight constraints, high mobility is required to escape local optimums. 

\textbf{Problem Setup:} The MSJD is defined as $\mathbb{E}[||x_{t+1} - x_t||^2]$. The target was a non-uniform financial risk model bounded by strict linear inequalities ($Aw \le b$), scaling from $D=10$ to $D=40$ dimensions.

\textbf{Graph Interpretation:} 
The x-axis represents the dimensionality of the target problem. The y-axis (log scale) represents the physical MSJD. A higher line indicates an algorithm that takes vastly larger steps through the physical volume of the polytope without crashing into walls.

\textbf{Observations:}
\begin{itemize}
    \item At $D=10$, WHSS achieved a jump distance of $1.29 \times 10^{-3}$, outperforming HRSS ($7.71 \times 10^{-4}$) and Dikin Walk ($1.22 \times 10^{-4}$).
    \item As dimensionality scales to $D=40$, the disparity widens. WHSS achieved $1.68 \times 10^{-4}$ compared to HRSS ($9.50 \times 10^{-5}$). This proves that WHSS maintains a higher mobility floor, escaping local traps where baselines are forced into small boundary steps.
\end{itemize}

\begin{figure}[H]
    \centering
    \includegraphics[width=0.8\textwidth]{msjd_plot.png}
    \caption{Mean Squared Jump Distance vs Dimensionality.}
\end{figure}

\section{Dimensional Scaling}
\textbf{Motivation:} Many MCMC algorithms perform well in low dimensions but suffer decay in efficiency (the curse of dimensionality) as the volume of geometric corners increases. We designed this experiment to test the dimensional scaling of these algorithms.

\textbf{Problem Setup:} The target was a skewed Gaussian (Condition Number = 1000) inside a strict bounding box $\mathcal{P} = \{x \in \mathbb{R}^D \mid -1 \le x_i \le 1\}$, evaluated across exponentially increasing scales: $D=10, 50, 100, 200,$ and $400$.

\textbf{Methodology and Metric Definition:}
All three algorithms were run for a fixed budget of $N_{\text{MCMC}} = 25{,}000$ post-warm-up MCMC steps. The reported metric is:
\[
\text{Efficiency} = \frac{\text{ESS}}{N_{\text{MCMC}} / 1000} \quad \text{(ESS per 1,000 MCMC steps)}
\]
This metric measures \emph{chain mixing quality}: how many statistically independent samples are produced per unit of MCMC work. This is the standard reporting convention in adaptive MCMC literature (e.g., the No-U-Turn Sampler, Hoffman et al., 2014; and HMC, Neal, 2011), analogous to how NUTS reports post-adaptation ESS rates while separately disclosing adaptation cost.

\textbf{Warm-Up Transparency:} WHSS employs a one-time geometric warm-up phase (Phases 1 and 2 of the PhaseManager) prior to MCMC sampling. This warm-up shoots exploratory rays across the polytope and fits a global warp matrix $L$ to the geometry. The warm-up budget scales as $\mathcal{O}(D)$ target evaluations (enforced internally as $\texttt{burn\_in\_samples} = 1000 \times D$). At $D=400$, this amounts to exactly $400{,}000$ target evaluations before a single MCMC step is taken. The baseline algorithms (HRSS, Dikin Walk) require \textbf{no} global warm-up and begin collecting MCMC samples from step 1.

It is critical to note that the label \emph{``post-warm-up''} does \textbf{not} imply that HRSS or Dikin Walk underwent any adaptation phase. Rather, it means we are measuring each algorithm's \emph{steady-state MCMC kernel efficiency}: for HRSS and Dikin Walk, that steady state is reached immediately at step 1 (no pre-computation exists); for WHSS, it is reached after the geometric warm-up completes. Both are then evaluated over the same $N_{\text{MCMC}} = 25{,}000$ MCMC steps. Furthermore, for robust measurement of steady-state mixing, the first 20\% (5{,}000 steps) of all MCMC chains were discarded as burn-in before calculating ESS. This universally eliminates any non-stationary transient drift from the initial state $x_0 = \mathbf{0}$. The warm-up cost for WHSS is \emph{not} included in the ESS/1k denominator above.

\textbf{Amortization Argument:} Because the WHSS warm-up cost is a fixed one-time investment and the MCMC mixing rate is thereafter constant, the warm-up fully amortizes in long-horizon workloads. Concretely, at $D=400$, an algorithm requiring target ESS $E$ needs:
\begin{align*}
\text{NFE}_{\text{HRSS}} &= 2083 \times E \\
\text{NFE}_{\text{WHSS}} &= 400{,}000 + 227 \times E
\end{align*}
WHSS breaks even with HRSS at $E^* \approx 215$ effective samples. For any scientific application requiring $E > 10{,}000$ ESS (the standard threshold for credible interval estimation), WHSS strictly dominates. As $E \to \infty$, the ratio $\text{NFE}_{\text{HRSS}} / \text{NFE}_{\text{WHSS}} \to 9.17\times$, reflecting the massive efficiency gap in the MCMC phase.

\textbf{Observations:}
\begin{table}[H]
\centering
\begin{tabular}{l|cc|cc|cc}
\toprule
\textbf{Dimension} & \multicolumn{2}{c|}{\textbf{WHSS (Ours)}} & \multicolumn{2}{c|}{\textbf{HRSS}} & \multicolumn{2}{c}{\textbf{Dikin Walk}} \\
 & ESS/1k & Time & ESS/1k & Time & ESS/1k & Time \\
\midrule
10D  & 36.44  & $<$ 1s & 30.65  & $<$ 1s & 17.09 & $<$ 1s \\
50D  & 7.43   & $<$ 1s & 1.99   & $<$ 3s & 1.05  & $<$ 5s \\
100D & 5.51   & $<$ 3s & 0.83   & $<$ 10s & 0.62 & $\sim$ 30s \\
200D & 5.32   & $<$ 5s & 0.58   & $\sim$ 30s & 0.84 & $\sim$ 2m \\
400D & \textbf{4.40}  & $<$ 8s & 0.48  & $\sim$ 2m & 0.44 & $\sim$ 13m \\
\bottomrule
\end{tabular}
\caption{Post-Warm-Up MCMC Mixing Efficiency and raw MCMC Wall-Clock Time (for $25{,}000$ samples). Warm-up cost for WHSS is $\mathcal{O}(D)$ and excluded from the MCMC time metric; see Warm-Up Transparency note above.}
\end{table}

\begin{figure}[H]
    \centering
    \includegraphics[width=0.8\textwidth]{scaling_plot.png}
    \caption{Dimensional Scaling: ESS per 1,000 NFE. WHSS maintains a steady efficiency floor while baseline algorithms suffer exponential decay in high dimensions.}
\end{figure}

At $10D$, the space is trivial and all algorithms perform similarly. By $100D$, HRSS undergoes exponential decay (falling below 1.0 ESS) because narrow geometric corners dominate the volume. By $400D$, the $\mathcal{O}(D^3)$ Hessian inversion inside Dikin Walk increases its computational time significantly: generating 25,000 samples for Dikin Walk required substantial computational time (the 400D matrix inversion loop is exceptionally slow), whereas WHSS finishes its post-warm-up MCMC phase in seconds due to its $\mathcal{O}(D^2)$ amortized step cost. WHSS maintains a dimension-independent floor of $\sim 4.40$ ESS/1k across all tested dimensions. This indicates that the one-time geometric warp insulates the MCMC chain from both statistical and computational collapse. \textbf{Crucially, this experiment demonstrates that if the underlying constraints are well-behaved (e.g., an isotropic hypercube or standard simplex), WHSS can scale to 400D and solve high target-skew (such as a 1000:1 Gaussian penalty). Its failure modes are primarily tied to highly skewed constraint boundaries, not high dimensionality or skewed target densities.}

\textbf{Algorithmic Trade-offs and Limitations (Disclaimer):} 
To maintain scientific transparency, we explicitly note the limitations of WHSS. WHSS requires a mandatory geometric warm-up phase of $\mathcal{O}(D)$ target evaluations to fit its affine warp matrix. At $D=400$, this warm-up requires $400{,}000$ function evaluations before a single MCMC step is generated. For extremely short-horizon tasks where only a few hundred samples are needed, this overhead is computationally prohibitive, and baseline algorithms like Hit-and-Run are preferred. WHSS is strictly designed for high-dimensional, long-horizon sampling where massive sample sizes are required, allowing the fixed one-time warm-up cost to amortize over the steady-state chain. Furthermore, if a target density is highly multimodal, the current global affine warp may fail to capture disparate modes, a known limitation of global preconditioners.

\section{Real-World Case Study: Risk-Constrained Portfolio Optimization}
\textbf{Motivation:} To prove WHSS is viable in production environments, we applied it to Institutional Portfolio Optimization. Complex risk models and strict regulatory boundaries render standard first-order algorithms ineffective because the constraint boundaries are non-differentiable. We included \textbf{Emcee} (the Affine-Invariant Ensemble Sampler) because it is a standard benchmark for gradient-free black-box sampling. We included Emcee to expose why naive boundary rejection fails in high dimensions.

\textbf{Problem Setup:} The model simulates a 30-asset portfolio bounded by strict regulatory constraints: a fully invested financial simplex ($\sum w_i = 1$), long-only bounds ($w_i \ge 0$), single-asset diversification caps ($w_i \le 0.15$), and a collective sector exposure limit ($\sum_{i=1}^5 w_i \le 0.30$). These constraints form a truncated simplex with narrow geometric corners. 

The target density evaluates a \textbf{Multi-Factor Financial Risk Model}. The covariance matrix ($\Sigma$) is constructed using 3 macro-economic market factors (with random loadings) plus idiosyncratic noise, mirroring industry standard Barra/Fama-French risk models. The sampler explores portfolios weighted by a Markowitz-style Risk-Aversion penalty: $f(w) \propto \exp(-\gamma \cdot w^T \Sigma w)$, where $\gamma=500.0$. This skewed target density runs tightly along the constraint walls, creating a highly constrained sampling space. We allocated a strict computational budget of exactly $N = 60{,}000$ samples per algorithm, averaged over 10 independent runs.

\textbf{Observations:}
\begin{table}[H]
\centering
\begin{tabular}{l|cc|c}
\toprule
\textbf{Algorithm} & \textbf{Min ESS} & \textbf{ESS/1k NFE} & \textbf{Domain Coverage (\%)} \\
\midrule
Emcee (Ensemble) & 669.1 & 11.15 & 79.3\% \\
HRSS & 168.1 & 2.80 & 100.0\% \\
Dikin Walk & 174.9 & 2.91 & 99.7\% \\
\textbf{WHSS (Ours)} & \textbf{210.6} & \textbf{3.51} & \textbf{100.0\%} \\
\bottomrule
\end{tabular}
\caption{Efficiency and Spatial Coverage on 30-Asset CVaR Target.}
\end{table}

At first glance, Emcee appears to perform well with 11.15 ESS/1k NFE. However, this is a misleading metric. Emcee handles constraints purely by rejection (assigning $-\infty$ density outside the boundary). In a high-dimensional simplex, the boundary volume is large. During our benchmarks, Emcee suffered an extremely low acceptance rate ($<10\%$), meaning over 90\% of its proposals were rejected. As a result, the ensemble collapsed into a localized subset of the polytope, ultimately exploring only $79.3\%$ of the valid domain. Its high ESS simply means it was mixing locally inside a trapped subspace. 

When comparing only the algorithms that successfully achieved $100\%$ domain coverage, WHSS outperforms the field. WHSS achieved $3.51$ ESS/1k NFE, compared to HRSS (2.80) and Dikin Walk (2.91).

\section{Real-World Case Study: Metabolic Flux Analysis}
\textbf{Motivation:} The sampling of biological metabolic networks represents a challenging constraint environment in computational biology. Biological polytopes are typically anisotropic, forming long, tight "needles" of feasible flux states. The standard framework for these spaces is \texttt{cobrapy}, which relies primarily on two algorithms for steady-state exploration: ACHR and CHRR. 

\subsection{Algorithmic Boundaries and Trade-offs (The Uniform Needle Trap)}
Before applying experimental data penalties, it is necessary to establish WHSS's algorithmic boundaries. We first evaluated WHSS on a purely Uniform target distribution ($f(v) \propto 1$) in the \texttt{e\_coli\_core} model constraint space.

\textbf{Result:} WHSS struggled to properly explore the uniform biological space, exhibiting a high Centroid Discrepancy (L2 norm distance of $140.49$ from the 200k empirical ground truth mean, compared to ACHR's $9.18$). 

\textbf{Mathematical Cause:} This is caused by a geometric vulnerability in WHSS's von Mises-Fisher (vMF) warm-up phase, which we term the ``Needle Trap.'' When placed in the center of a highly skewed biological polytope, isotropic vMF rays overwhelmingly collide with the side-walls of the needle. Almost zero rays traverse the long axis. Consequently, the empirical covariance matrix generated from these boundary anchors severely underestimates the variance along the long axis. The $L$-matrix incorrectly warps the space into a tight box, trapping the slice sampler in a false-stationary local optimum.

\textbf{Conclusion:} For pure uniform exploration of highly constrained biological spaces, WHSS is ineffective. The current State-of-the-Art, CHRR (Coordinate Hit-and-Run with Rounding)---which actively calculates a Maximum Volume Ellipsoid to round the entire space---remains the standard approach.

\subsection{The Non-Uniform Breakthrough: Transcriptomics Integration}
While CHRR solves the uniform sampling problem, modern computational biology increasingly requires sampling \emph{non-uniform}, data-penalized target densities. 

\textbf{Problem Setup:} To simulate the integration of RNA-seq or 13C-MFA isotope data, we apply a Gaussian penalty ($\sigma=1.0$) centered on the FBA optimum of the 24-dimensional \texttt{e\_coli\_core} model: $f(v) \propto \exp(-\|v - v_{\text{opt}}\|^2 / 2\sigma^2)$. Total budget $N=25{,}000$; 50\% burn-in discarded before ESS calculation.

\textbf{The SOTA Compromise:} Because tools like \texttt{cobrapy} only provide Uniform samplers (ACHR/CHRR) and standard HMC fails against strict inequality boundaries, a widely discussed theoretical approach is to sample uniformly using CHRR, and then re-weight the samples using Sampling Importance Resampling (SIR) to approximate the non-uniform density (e.g., following frameworks by Antoniewicz, 2015; Haraldsd\'ottir et al., 2017). We tested this protocol against native WHSS.

\textbf{Experimental Rationale:} The 24D model is used specifically because it is the smallest published model with a complete null space, allowing us to generate a computationally verifiable ground truth. The weight collapse observed at 24D serves as a lower bound. Weight degeneracy in SIR grows exponentially with dimensionality; extrapolating to genome-scale (542D), SIR is intractable.

\begin{table}[H]
\centering
\begin{tabular}{l|ccc|c}
\toprule
\textbf{Algorithm} & \textbf{Time (s)} & \textbf{Weight Retention} & \textbf{Min ESS} & \textbf{ESS/1k NFE} \\
\midrule
ACHR (Uniform+SIR) & 0.64 & 0.01\% $^\dagger$ & 1.0 & 0.04 \\
CHRR (Uniform+SIR) & 2.92 & 0.01\% $^\dagger$ & 1.0 & 0.04 \\
\textbf{WHSS (Native)} & 4.48 & \textbf{100.0\%} & \textbf{102.8} & \textbf{4.11} \\
\bottomrule
\end{tabular}
\caption{24D Non-Uniform Transcriptomics Integration. $^\dagger$ Weight Retention $<5\%$ renders SIR-based distributions mathematically inefficient regardless of MCMC ESS (Kish, 1965).}
\end{table}

\textbf{Discussion:} The results reveal that the tested protocol (CHRR + SIR) is ineffective in high dimensions. Because the Gaussian penalty is tightly concentrated, the vast majority of uniformly scattered points receive negligible weight, collapsing the effective sample size to $\approx 1.0$ out of 25,000 samples ($0.01\%$ retention). 

\textbf{Mathematical Definition of Weight Retention:}
In Sampling Importance Resampling (SIR), samples drawn from a uniform proposal distribution $Q(x)$ must be re-weighted to match the true non-uniform target density $P(x)$ using importance weights $w_i \propto P(x_i)/Q(x_i)$. The statistical efficiency of the resulting ensemble is governed by the Kish Effective Sample Size (Kish, 1965):
\[
\text{ESS}_{IS} = \frac{\left(\sum_{i=1}^N w_i\right)^2}{\sum_{i=1}^N w_i^2}
\]
\emph{Weight Retention} is defined as the percentage of the raw sample budget $N$ that survives this re-weighting process: $\text{Retention} = (\text{ESS}_{IS} / N) \times 100\%$. 
A retention of $0.01\%$ means the target density is so concentrated relative to the uniform space that out of $25{,}000$ generated samples, the weights collapsed onto a single point ($\text{ESS}_{IS} \approx 1.0$). The remaining $24{,}999$ samples were effectively assigned a weight of zero and discarded, resulting in a statistically inefficient sample.

Conversely, the Gaussian penalty resolves WHSS's Needle Trap. While the underlying biological polytope remains skewed, the penalty function acts as a bounding region, removing the un-explorable ends of the needle. The \emph{effective} target space becomes tightly concentrated and isotropic around the data. This allows WHSS's vMF ray-casting to map the relevant probability mass efficiently. Because WHSS evaluates the non-uniform density natively at every step, it requires zero re-weighting ($100\%$ retention) and generates high MCMC efficiency (4.11 ESS/1k NFE) in an environment where baseline algorithms struggle.

\end{document}
